% Part: normal-modal-logic
% Chapter: axioms-systems
% Section: derived-rules

\documentclass[../../../include/open-logic-section]{subfiles}

\begin{document}

\olfileid{nml}{prf}{der}

\olsection{వ్యుత్పన్న నియమాలు}

\tetoken{వ్యుత్పత్తులను}{!!{derivation}s} కనుగొని రాయడం కష్టసాధ్యం,
శ్రమతో కూడినది, పునరావృతమయ్యేది కూడా. ఉదాహరణకు, $!A \lif !B$ నుంచి
$\Box !A \lif \Box !B$కు వెళ్లాలని, అంటే నియమం~\RKను
వర్తింపజేయాలని, మనకు తరచుగా అవసరమవుతుంది. అందుకు ముందుగా \Necను
వర్తింపజేసి, తరువాత~\Ax{K} యొక్క తగిన ప్రతిస్థాపన నిదర్శనాన్ని
నమోదు చేసి, చివరగా~\MPను వర్తింపజేయాలి. $!A \lif !B$, $!B \lif !C$
నుంచి $!A \lif !C$కు వెళ్లాలంటే, ఈ (పొడవైన) సర్వసత్య నిదర్శనాన్ని
నమోదు చేయాలి:
\[
(!A \lif !B) \lif ((!B \lif !C) \lif (!A \lif !C))
\]
దీనికి \MPను రెండుసార్లు వర్తింపజేయాలి. మనకు తెలిసిన సమానార్థక
సూత్రంతో ఒక \tetoken{ఉపసూత్రాన్ని}{!!{formula}} భర్తీ చేయాలని కూడా
తరచుగా అనుకుంటాం; ఉదాహరణకు,
\iftag{prvDiamond}{$\Diamond
  !A$ స్థానంలో $\lnot\Box\lnot !A$}{$\Box !A$ స్థానంలో
  $\lnot\Diamond\lnot !A$}, లేదా $\lnot\lnot !A$ స్థానంలో $!A$.
అందువల్ల పూర్తి \tetoken{వ్యుత్పత్తిని}{!!{derivation}} రాయడంకంటే,
మధ్యంతర దశలు ఎందుకు \tetoken{వ్యుత్పాదించదగినవో}{!!{derivable}}
సూచించడం సౌకర్యంగా ఉంటుంది. ఇందుకోసం
\tetoken{వ్యుత్పాద్యత}{!!{derivability}}కు సంబంధించిన కొన్ని
వాస్తవాలను సమీకరిద్దాం.

\begin{prop}
  $\Log{K} \Proves !A_1$, \dots, $\Log{K} \Proves !A_n$ కాగా,
  ప్రతిజ్ఞావాక్య తర్కం ప్రకారం $!A_1$, \dots, $!A_n$ నుంచి $!B$
  అనుసరిస్తే, $\Log{K} \Proves !B$.
\end{prop}

\begin{proof}
  ప్రతిజ్ఞావాక్య తర్కం ప్రకారం $!A_1$, \dots,~$!A_n$ నుంచి $!B$
  అనుసరిస్తే,
  \[
  !A_1 \lif (!A_2 \lif \cdots (!A_n \lif !B)\dots)
  \]
  ఒక సర్వసత్య ప్రతిస్థాపన నిదర్శనం. \MPను $n$ సార్లు వర్తింపజేస్తే
  $!B$కు ఒక \tetoken{వ్యుత్పత్తి}{!!{derivation}} లభిస్తుంది.
\end{proof}

ఈ ప్రతిపాదనను వాడినప్పుడు~\PLతో సూచిస్తాం.

\begin{prop}
  $\Log{K} \Proves !A_1 \lif (!A_2 \lif \cdots (!A_{n-1} \lif
  !A_n)\dots)$ అయితే $\Log{K} \Proves \Box !A_1 \lif (\Box !A_2 \lif
  \cdots (\Box !A_{n-1} \lif \Box !A_n)\dots)$.
\end{prop}

\begin{proof}
  \olref[nor]{prop:rk} నిరూపణలో మాదిరిగానే $n$పై ఆగమనం ద్వారా.
\end{proof}

ఈ ప్రతిపాదనను వాడినప్పుడు~\RKతో సూచిస్తాం. ఈ ఫలితాలు
\tetoken{వ్యుత్పాద్యత}{!!{derivability}} ఫలితాలను మరింత సులభంగా
స్థాపించడానికి ఎలా తోడ్పడతాయో చూద్దాం.

\begin{prop}
  $\Log{K} \Proves (\Box!A \land \Box!B) \lif \Box (!A \land !B)$
\end{prop}

\begin{proof}
  \begin{derivation}
    1. & $\Log{K} \Proves !A \lif (!B \lif (!A \land !B))$ & \Taut \\
    2. & $\Log{K} \Proves \Box!A \lif (\Box !B \lif \Box(!A \land !B)))$ & \RK, 1\\
    3. & $\Log{K} \Proves (\Box!A \land \Box!B) \lif \Box(!A \land !B)$ & \PL, 2
  \end{derivation}
\end{proof}

\begin{prop}\ollabel{prop:rewriting}
  $\Log{K} \Proves !A \liff !B$ మరియు $\Log{K} \Proves
  \Subst{!C}{!A}{q}$ అయితే $\Log{K} \Proves \Subst{!C}{!B}{q}$.
  \sourcecorrection{OLTENMLAXSDER-001}{స్థిర మూలంలోని ఈ
    తీర్మానంలో రెండవ ప్రతిస్థాపనకు $B$ అని మాత్రమే ఉంది;
    తరువాతి అభ్యాసం, ఇక్కడి పరికల్పన రెండూ చూపినట్లుగా
    సూత్ర మెటాచరం $!B$ను పునరుద్ధరించాం.}
\end{prop}

\begin{proof}
  వ్యాయామం.
\end{proof}

\begin{prob}
  $!C$ సంక్లిష్టతపై ఆగమనంతో, $\Log{K} \Proves !A \liff !B$ అయితే
  $\Log{K} \Proves \Subst{!C}{!A}{q} \liff \Subst{!C}{!B}{q}$ అని
  నిరూపించి, \olref[nml][prf][der]{prop:rewriting}ను నిరూపించండి.
\end{prob}

ముఖ్యంగా $\Diamond$ను $\Box$గా (లేదా విలోమంగా) మార్చాలన్నా,
\tetoken{ఒక సూత్రం}{!!a{formula}} లోపలి ద్వినిషేధాలను తొలగించాలన్నా
ఈ ప్రతిపాదన ఉపయోగపడుతుంది. ఇకపై \olref{prop:rewriting}ను
వర్తింపజేసి $!C(!A)$ అనే \tetoken{సూత్రాన్ని}{!!a{formula}}
$!C(!B)$గా తిరిగి రాసినప్పుడు,
``$!A$ స్థానంలో $!B$'' అని సూచిస్తాం.
\sourcecorrection{OLTENMLAXSDER-002}{స్థిర మూలంలో సాధారణ
  సంక్షిప్త సూచన ``$!A$ for $!B$'' అని ఉంది; కానీ చూపిన
  దిశ $!C(!A)$ నుంచి $!C(!B)$కు, తరువాతి ఉదాహరణలోనూ
  పాత పదం స్థానంలో కొత్త పదమే వాడారు. కనుక దిశను
  ``$!A$ స్థానంలో $!B$''గా స్పష్టం చేశాం.}
మరో మాటలో, ``$!A$ స్థానంలో $!B$'' అనేది దీని సంక్షిప్త రూపం:
  \begin{derivation}
    & $\Proves !C(!A)$ \\
    & $\Proves !A \liff !B$\\
    & $\Proves !C(!B)$ & \olref{prop:rewriting} ప్రకారం
  \end{derivation}
ఉదాహరణకు:

\begin{prop}
  $\Log{K} \Proves \lnot\Box p \lif \Diamond \lnot p$
\end{prop}

\begin{proof}
  \iftag{prvDiamond}{% Diamond is primitive
  \begin{derivation}
    1. & $\Log{K} \Proves \Diamond \lnot p \liff \lnot\Box\lnot\lnot p$ &
    \Dual\\
    2. & $\Log{K} \Proves \lnot \Box \lnot\lnot p \lif \Diamond \lnot p$ & \PL, 1\\
    3. & $\Log{K} \Proves \lnot\Box p \lif \Diamond \lnot p$ & $\lnot\lnot p$ స్థానంలో $p$\\
  \end{derivation}
  }{% Diamond is defined
  \begin{derivation}
    1. & $\Log{K} \Proves \lnot \Box \lnot\lnot p \lif \Diamond \lnot p$ & \Taut\\
    2. & $\Log{K} \Proves \lnot\Box p \lif \Diamond \lnot p$ & $\lnot\lnot p$ స్థానంలో $p$\\
  \end{derivation}
  $1$వ పంక్తిలోని సూత్రం, $q$ స్థానంలో $\Box\lnot\lnot p$ను
  ప్రతిస్థాపిస్తే వచ్చే $\lnot q \lif \lnot q$ నిదర్శనం.
  $\Diamond\lnot p$, $\lnot\Box\lnot\lnot p$ రెండూ ఒకటే;
  ఎందుకంటే $\Diamond$ను $\lnot\Box\lnot$గా నిర్వచించాం.}
\end{proof}

పై \tetoken{వ్యుత్పత్తి}{!!{derivation}}లో చివరి దశ
``$\lnot\lnot p$ స్థానంలో $p$'' అనేది దీని సంక్షిప్త రూపం:
  \begin{derivation}
    & $\Log{K} \Proves \lnot \Box \lnot\lnot p \lif \Diamond \lnot
    p$\\
    & $\Log{K} \Proves \lnot\lnot p \liff p$ & \Taut\\
    & $\Log{K} \Proves \lnot\Box p \lif \Diamond \lnot p$ & \olref{prop:rewriting} ప్రకారం
  \end{derivation}
ఇక్కడ \olref{prop:rewriting}లోని $!C(q)$, $!A$, $!B$ల స్థానాల్లో
వరుసగా $\lnot \Box q \lif \Diamond \lnot p$,
$\lnot\lnot p$, $p$ ఉన్నాయి.

\tetoken{ఒక సూత్రంలో}{!!a{formula}} $\lnot\Diamond !A$ అనే
\tetoken{ఉపసూత్రం}{!!{formula}} ఉంటే,
$\Log{K} \Proves \lnot\Diamond !A \liff \Box\lnot !A$ కనుక,
\olref{prop:rewriting}ను వాడి దాని స్థానంలో $\Box\lnot !A$ను
రాయవచ్చు. ఇటువంటి భర్తీలను ``$\lnot\Diamond$ స్థానంలో
$\Box\lnot$'' అని సంక్షిప్తంగా సూచిస్తాం.

కింది ప్రతిపాదన, \tetoken{వ్యుత్పాద్యత}{!!{derivability}} ఫలితాలను
సూత్రపద్ధతిలో స్థాపించవచ్చని సమర్థిస్తుంది. ఉదాహరణకు, పూర్వ
ప్రతిపాదన కేవలం $\Log{K} \Proves \lnot\Box p \lif
\Diamond \lnot p$ను మాత్రమే కాక, ఏ $!A$కైనా
$\Log{K} \Proves \lnot\Box !A \lif \Diamond \lnot !A$ను
స్థాపిస్తుంది.

\begin{prop}
  $!A$, $!B$కు ప్రతిస్థాపన నిదర్శనమై, $\Log{K} \Proves !B$ అయితే,
  $\Log{K} \Proves !A$.
\end{prop}

\begin{proof}
  $!B$కు ఉన్న \tetoken{వ్యుత్పత్తి}{!!{derivation}} పొడవుపై
  ఆగమనం చేసి, మొత్తం \tetoken{వ్యుత్పత్తికి}{!!{derivation}}
  ప్రతిస్థాపన వర్తింపజేసినప్పుడు కూడా సరైన
  \tetoken{వ్యుత్పత్తి}{!!{derivation}} లభిస్తుందని తనిఖీ చేయడం
  శ్రమతో కూడినదైనా సాధారణమైన పని. ప్రత్యేకంగా, సర్వసత్య నిదర్శనాల
  ప్రతిస్థాపన నిదర్శనాలు మళ్లీ సర్వసత్య నిదర్శనాలే;
  \iftag{prvDiamond}{\Dual{} మరియు~}{}\Ax{K} నిదర్శనాల
  ప్రతిస్థాపన నిదర్శనాలు మళ్లీ
  \iftag{prvDiamond}{\Dual{} మరియు~}{}\Ax{K} నిదర్శనాలే; అలాగే
  పూర్వపక్షం(లు), తీర్మానం రెండింటిలో
  \tetoken{ప్రతిజ్ఞావాక్య చరాల}{!!{propositional variable}s}
  స్థానంలో \tetoken{సూత్రాలను}{!!{formula}s} ప్రతిస్థాపించినా
  \MP{}, \Nec{} అనువర్తనాలు సరైనవిగానే ఉంటాయి.
\end{proof}


\end{document}
